\documentclass[14pt,a4paper,oneside]{extbook}
\usepackage[utf8]{vietnam}
\usepackage{amsmath,amssymb,mathrsfs,fancyhdr,enumerate,multirow,makecell}
\usepackage{tkz-euclide,tikz-3dplot,tikz,tkz-tab}
\usetikzlibrary{arrows.meta,calc,intersections,angles,quotes,backgrounds,shapes.geometric,patterns,through,shadows,positioning}
\usepackage{yhmath}
\usepackage[top=0.7cm, bottom=0.7cm, left=1.0cm, right=1.0cm, headheight=16pt, headsep=3mm, footskip=5mm]{geometry}
\usepackage[hidelinks,unicode]{hyperref}
\usepackage{qrcode}
\usepackage[many]{tcolorbox}
\usepackage{fontawesome}

\pagestyle{fancy}
\fancyhf{}
\fancyhead[L]{\footnotesize\bfseries\color{teal!85!black}\faBook~SỔ TAY KIẾN THỨC TOÁN 4}
\fancyhead[R]{\footnotesize\bfseries\color{blue!85!black}Lớp Toán Thầy Lĩnh -- loptoanthaylinh.vn}
\fancyfoot[C]{\footnotesize\bfseries\color{gray!80!black}-- Sổ tay kiến thức: Bài 38. Nhân với số có một chữ số --}
\renewcommand{\headrulewidth}{0.4pt}

\begin{document}

% --- BANNER TIÊU ĐỀ BÀI HỌC ---
\noindent
\begin{tcolorbox}[
  enhanced,
  colback=teal!8!white,
  colframe=teal!80!black,
  arc=3mm, boxrule=1pt,
  left=3mm, right=3mm, top=2mm, bottom=2mm,
  halign=center
]
  {\Large\bfseries\color{teal!90!black}\faBookmark\ BÀI 38: NHÂN VỚI SỐ CÓ MỘT CHỮ SỐ}\\[1.5pt]
  {\footnotesize\itshape\color{black!75} Kĩ thuật đặt tính thẳng cột, quy tắc nhân từ phải sang trái và xử lí nhớ sang hàng liền trước trên 1 trang A4}
\end{tcolorbox}

\vspace{2mm}

% --- KHỐI 1 & 2: LÝ THUYẾT & KĨ NĂNG ---
\noindent
\begin{minipage}[t]{0.49\textwidth}
  \begin{tcolorbox}[
    enhanced,
    title={\footnotesize\bfseries\color{white}\faPencil\ 1. Đặt tính \& Thứ tự thực hiện},
    colback=blue!2!white,
    colframe=blue!75!black,
    arc=2.5mm, boxrule=0.8pt, top=1.8mm, bottom=1.8mm, left=2mm, right=2mm
  ]
    \footnotesize
    $\bullet$ \textbf{Đặt tính}:
    \par
    --- Viết thừa số thứ nhất ở trên.\\
    --- Viết thừa số thứ hai (có 1 chữ số) ở dưới, \textbf{thẳng cột hàng đơn vị}.\\
    --- Viết dấu $\times$ bên trái, kẻ dấu gạch ngang dưới.
    $\bullet$ \textbf{Thứ tự tính}: Thực hiện nhân \textbf{từ phải sang trái} (từ hàng đơn vị đến chục, trăm, nghìn...).
  \end{tcolorbox}
\end{minipage}%
\hfill
\begin{minipage}[t]{0.49\textwidth}
  \begin{tcolorbox}[
    enhanced,
    title={\footnotesize\bfseries\color{white}\faLightbulbO\ 2. Kĩ thuật nhân có nhớ},
    colback=teal!2!white,
    colframe=teal!75!black,
    arc=2.5mm, boxrule=0.8pt, top=1.8mm, bottom=1.8mm, left=2mm, right=2mm
  ]
    \footnotesize
    $\bullet$ Nếu tích của một hàng $\ge 10$:
    \par
    --- Viết chữ số hàng đơn vị vào kết quả.\\
    --- \textbf{Nhớ chữ số hàng chục} sang hàng liền kề bên trái.
    $\bullet$ \textbf{Quy tắc vàng}: Lượt nhân tiếp theo phải \textbf{nhân trước rồi mới cộng số nhớ vào tích}.\\
    $\bullet$ Lượt nhân cuối cùng: Viết toàn bộ kết quả.
  \end{tcolorbox}
\end{minipage}

\vspace{2mm}

% --- KHỐI 3: SƠ ĐỒ TƯ DUY TIKZ ---
\noindent
\begin{tcolorbox}[
  enhanced,
  title={\footnotesize\bfseries\color{white}\faSitemap\ SƠ ĐỒ TƯ DUY: QUY TRÌNH NHÂN CÓ NHỚ $14\,253 \times 4$ TRỰC QUAN},
  colback=yellow!2!white,
  colframe=orange!85!black,
  arc=2.5mm, boxrule=0.8pt, top=1.5mm, bottom=1.5mm, left=1mm, right=1mm,
  halign=center
]
  \begin{tikzpicture}[
    scale=0.92,
    box/.style={rectangle, rounded corners=2mm, draw=teal!80!black, fill=teal!10, thick, inner sep=2.5pt, font=\scriptsize\bfseries, align=center},
    stepb/.style={rectangle, rounded corners=2mm, draw=blue!75!black, fill=blue!8, thick, inner sep=2pt, font=\scriptsize\bfseries, align=left, text width=9.2cm},
    arr/.style={-{Stealth[scale=0.85]}, thick, draw=black!70, shorten >= 2.5pt, shorten <= 2.5pt}
  ]
    % Cột phép tính bên trái
    \node[box, fill=white, draw=purple!80!black, text width=4.5cm, inner sep=4pt] (calc) at (2.2, 0) {
      \large
      \begin{tabular}{c@{\,}r@{\,}r@{\,}r@{\,}r@{\,}r}
        & $1$ & $4$ & $2$ & $5$ & $3$ \\
        $\times$ & & & & & $4$ \\
        \hline
        & $\mathbf{5}$ & $\mathbf{7}$ & $\mathbf{0}$ & $\mathbf{1}$ & $\mathbf{2}$
      \end{tabular}
    };

    % Các bước bên phải
    \node[stepb, draw=teal!80!black, fill=teal!8] (s1) at (9.8, 1.25) {\textbf{B1 (Đơn vị):} $4 \times 3 = 12 \rightarrow$ Viết $\mathbf{2}$, nhớ $1$};
    \node[stepb, draw=orange!80!black, fill=orange!10] (s2) at (9.8, 0.42) {\textbf{B2 (Hàng chục):} $4 \times 5 = 20$, thêm $1$ bằng $21 \rightarrow$ Viết $\mathbf{1}$, nhớ $2$};
    \node[stepb, draw=blue!75!black, fill=blue!8] (s3) at (9.8, -0.42) {\textbf{B3 (Hàng trăm):} $4 \times 2 = 8$, thêm $2$ bằng $10 \rightarrow$ Viết $\mathbf{0}$, nhớ $1$};
    \node[stepb, draw=purple!80!black, fill=purple!10] (s4) at (9.8, -1.25) {\textbf{B4 (Nghìn \& Chục nghìn):} $4 \times 4 + 1 = 17$ (viết $\mathbf{7}$, nhớ $1$); $4 \times 1 + 1 = \mathbf{5}$};

    \draw[arr] (calc.east) -- (s2.west);
  \end{tikzpicture}
\end{tcolorbox}

\vspace{2mm}

% --- KHỐI 4: VÍ DỤ MINH HỌA & BẪY TRẮC NGHIỆM ---
\noindent
\begin{minipage}[t]{0.53\textwidth}
  \begin{tcolorbox}[
    enhanced,
    title={\footnotesize\bfseries\color{white}\faCheckCircle\ Ví dụ thực hành chuẩn mực},
    colback=green!2!white,
    colframe=green!60!black,
    arc=2.5mm, boxrule=0.8pt, top=1.8mm, bottom=1.8mm, left=2.2mm, right=2.2mm
  ]
    \footnotesize
    \textbf{Ví dụ 1}: Đặt tính rồi tính $26\,148 \times 3$:\\
    --- $3 \times 8 = 24$, viết $4$ nhớ $2$.\\
    --- $3 \times 4 = 12$, thêm $2$ bằng $14$, viết $4$ nhớ $1$.\\
    --- $3 \times 1 = 3$, thêm $1$ bằng $4$, viết $4$.\\
    --- $3 \times 6 = 18$, viết $8$ nhớ $1$.\\
    --- $3 \times 2 = 6$, thêm $1$ bằng $7$, viết $7$.\\
    $\Rightarrow \mathbf{26\,148 \times 3 = 78\,444}$.
  \end{tcolorbox}
\end{minipage}%
\hfill
\begin{minipage}[t]{0.45\textwidth}
  \begin{tcolorbox}[
    enhanced,
    title={\footnotesize\bfseries\color{white}\faExclamationTriangle\ Mẹo tránh bẫy trắc nghiệm},
    colback=red!2!white,
    colframe=red!70!black,
    arc=2.5mm, boxrule=0.8pt, top=1.8mm, bottom=1.8mm, left=2.2mm, right=2.2mm
  ]
    \footnotesize
    $\bullet$ \textbf{Bẫy cộng nhớ trước khi nhân}:\\
    Học sinh hay lấy $(5 + 1) \times 4 = 24$ là \textbf{SAI}! Phải nhân trước: $4 \times 5 = 20$, rồi mới $+ 1 = 21$.\\
    $\bullet$ \textbf{Bẫy quên nhớ}: Khi nhân số có nhiều chữ số liên tiếp có nhớ, hãy chấm nhẹ lên đầu chữ số để không quên!
  \end{tcolorbox}
\end{minipage}

\end{document}
